\nuevaReceta{Caldo de Pescado}{1,5 l}{50 min}{receta:caldo_pescado}

    \begin{minipage}[t]{0.35\textwidth}
        \section*{Ingredientes}
        \begin{itemize}[leftmargin=*]
            \item 800 g de espinas y cabeza de pescado blanco, aprox.
            \item 1 cebolla
            \item 1 puerro
            \item 1 zanahoria
            \item 1 tomate maduro
            \item 2 dientes de ajo
            \item 1 hoja de laurel
            \item perejil fresco
            \item aceite de oliva
            \item 1,5 litros de agua
            \item sal
        \end{itemize}
    \end{minipage}
    \hfill
    \begin{minipage}[t]{0.6\textwidth}
        \section*{Preparación}
        \begin{enumerate}
            \item \textbf{Limpiar el pescado:} Lavar bien las espinas y la cabeza, retirando las agallas y cualquier resto de sangre para evitar sabores amargos.
            \item \textbf{Preparar las verduras:} Pelar y cortar la cebolla, el puerro, la zanahoria, el tomate y los ajos en trozos grandes.
            \item \textbf{Sofrito:} Calentar un poco de aceite en una olla y rehogar las verduras durante 5-8 minutos.
            \item \textbf{Añadir el pescado:} Incorporar las espinas y la cabeza. Cocinar un par de minutos junto con las verduras.
            \item \textbf{Cocción:} Añadir el agua, el laurel y el perejil. Cuando empiece a hervir, retirar la espuma de la superficie.
            \item \textbf{Hervor suave:} Cocer a fuego medio-bajo durante 25-30 minutos. No prolongar demasiado la cocción para que el caldo no quede amargo.
            \item \textbf{Colar:} Pasar el caldo por un colador fino, presionando ligeramente los ingredientes, y ajustar de sal.
        \end{enumerate}
    \end{minipage}

    \vspace{1em}
    \begin{tcolorbox}[colback=orange!5, colframe=orange!50, title=Consejo, size=small]
        Puedes congelar el caldo en porciones y utilizarlo como base para arroces, sopas, guisos o salsas de pescado.
    \end{tcolorbox}
\end{receta}
